\documentclass[11pt]{article}
% \usepackage[utf8]{inputenc} % Required for inputting international characters
\usepackage[T1]{fontenc} % Output font encoding for international characters
% % \usepackage{stix} % Use the STIX fonts
% % gives options for document layout
\usepackage{geometry}
\geometry{
    a4paper,
    top=20mm,
    right=23mm,
    left=23mm,
    }
% % \usepackage{mathtools}  % amsmath with extensions
% % \usepackage{amsfonts}
% \usepackage{amsmath}
% \usepackage{amssymb}
% \usepackage{bbold}
% % disables automatic indent at beginning of paragraph
\usepackage{parskip}
% % fancy headers
\usepackage{fancyhdr}
% % nice tables
% \usepackage{booktabs}
% % enable .png for \includegraphcis
\usepackage{graphicx}
% % subfigures
% % \usepackage[font=scriptsize,labelfont=bf,justification=justified]{caption}
% % \usepackage{subcaption}
% % \renewcommand{\thetable}{\alph{table}}
% % enable H to stop figures from floating
% % tables
% % \usepackage{tabularx}
% % make \left \right obsolete
% % \usepackage{nath}
% % make clickable links
% % \usepackage[linktocpage, hidelinks]{hyperref}
% \usepackage{hyperref}
% % all the good stuff
\usepackage{physics}
% % \usepackage{cite}
% % \usepackage[font=normal,labelfont=bf]{caption}
% \usepackage{xcolor}
% \hypersetup{
%     colorlinks,
%     linkcolor={black},
%     citecolor={black},
%     urlcolor={blue!80!black}
% }
% \usepackage[backend=biber,style=numeric,sorting=none,citestyle=ieee]{biblatex}


% \title{Volterra Equations}
% \author{Lucas Tepper$^1$, Cihan Ayaz$^1$ \vspace{0.2 cm} \\
% \emph{$^1$Department of Physics, Freie Universität Berlin, Germany}}
% \makedate
\begin{document}
% \maketitle

\section{Introduction}
\subsection{Markovian Embedding}
In this short report, I test the Volterra method that evaluates
the running integral $G(t) := \Gamma(t)$ on the half-steps $t_n = n\Delta t + \Delta t/2$,
using the trapezoidal rule.
\begin{equation}
    G_{i - \frac{1}{2}} = \frac{\Delta t}{2} \Gamma_0
        + \Delta t \int_{j=0}^{i - 1} \Gamma_j
        + \frac{\Delta t}{2} \Gamma_{i}
    \label{eq:discretization}
\end{equation}
We consider the Volterra equation used before
\begin{equation}
    m \left( C^{vv}(t) - C^{vv}(0) \right) =
        C^{xdU}(t) - C^{xdU}(0) - \int_{0}^{t} \dd s C^{vv}(t - j) \Gamma(j).
\end{equation}
Using equation~\ref{eq:discretization} to discretize the integral and assuming
$C^{vv}_{i + \frac{1}{2}} = \frac{1}{2} \left(C^{vv}_{i} + C^{vv}_{i + 1} \right)$, we obtain
\begin{equation}
    G_{i - \frac{1}{2}} = \frac{2}{\Delta t (C^{vv}_0 + C^{vv}_1)} \left(
        m \left( C^{vv}(0) - C^{vv}(t) \right) + C^{xdu}(t) - C^{xdu}(0)
        - \frac{\Delta t}{2} \sum_{j = 1}^{i - 1} (C^{vv}_j + C^{vv}_{j + 1}) G_{i - j - \frac{1}{2}}
        \right).
    \label{eq:volt_half}
\end{equation}
When setting $n = 0$ into the half-step method in equation~\ref{eq:volt_full}, we obtain
$G_{- \frac{1}{2}} = \int_{-\frac{1}{2}}^{0} = 0$. Given that we assumed the $\Gamma(t)$
to be symetric, this result is non-sensical. I therefore assume that the first element of
the numerical result is $G(0)$. I have verified that not using this assumption is unable to reproduce
the reference kernels in the results section, but omit to show it.
We want to compare it to the method evaluated on the full steps $t_n = n\Delta t$,
\begin{equation}
    G_{i} = \frac{2}{\Delta t C^{vv}_0} \left(
        m \left( C^{vv}(0) - C^{vv}(t) \right) + C^{xdu}(t) - C^{xdu}(0)
        - \Delta t \sum_{j = 1}^{i - 1} C^{vv}_j G_{i - j}
        \right).
    \label{eq:volt_full}
\end{equation}


\section{Results}
First, we test both methods for a Langevin simulation with a single delta
kernel $\Gamma(t) = 2 \gamma \delta(t)$ with $m = 31.4$ u, $\gamma = 5000$ u/ps
and simulate for 5e8 steps with a time step of 0.001 ps.
\begin{figure}[H]
    \centering
    \includegraphics{/net/data02/greedisgod/paper1/figures/halfstep_delta_kernel.pdf}
    \caption{We extract a memory kernel from equations 3 and 4.
    Without a potential, the $C^{xdu}$ terms are zero.
    We have a delta kernel with $\gamma = 5000$ u/ps.
    We shift time by dt to visualize $\Gamma(0)$.}
\end{figure}
Next, we test the same thing in the presence of an exponential kernel
$\Gamma(t) = \sum_i \frac{\gamma_i}{\tau_i} \exp(-\frac{t}{\tau_i})$ and a delta kernel. The
The simulation setup is the same, except for the two-mode exponential kernel with
$\gamma_i$ of [5000.0, 2000.0] and $\tau_i$ = [0.1, 1.0] ps.
\begin{figure}[H]
    \centering
    \includegraphics{/net/data02/greedisgod/paper1/figures/halfstep_delta_plus_exp.pdf}
    \caption{We extract a memory kernel from equations 3 and 4
    without a potential, the $C^{xdu}$ terms are zero.
    We have a delta kernel with $\gamma = 5000$ u/ps and an exponential kernel with
    $\gamma_i$ of [5000.0, 2000.0] and $\tau_i$ = [0.1, 1.0] ps
    We shift time by dt to visualize $\Gamma(0)$.}
\end{figure}
Next, we do the same tests in the potential landscape of Ala9.
\begin{figure}[H]
    \centering
    \includegraphics{/net/data02/greedisgod/paper1/figures/halfstep_ala9_dtelta.pdf}
    \caption{We extract a memory kernel from equations 3 and 4 in the ala9 potential.
    We have a delta kernel with $\gamma = 5000$ u/ps.
    We shift time by dt to visualize $\Gamma(0)$.}
\end{figure}
\begin{figure}[H]
    \centering
    \includegraphics{/net/data02/greedisgod/paper1/figures/halfstep_ala9_delta_plus_exp.pdf}
    \caption{We extract a memory kernel from equations 3 and 4 in the ala9 potential.
    We have a delta kernel with $\gamma = 5000$ u/ps and an exponential kernel with
    $\gamma_i$ of [5000.0, 2000.0] and $\tau_i$ = [0.1, 1.0] ps.
    We shift time by dt to visualize $\Gamma(0)$.}
\end{figure}
As a final test, we test the method in the presence of only an exponential kernel.
\begin{figure}[H]
    \centering
    \includegraphics{/net/data02/greedisgod/paper1/figures/halfstep_ala9_exp.pdf}
    \caption{We extract a memory kernel from equations 3 and 4 in the ala9 potential.
    We have an exponential kernel with
    $\gamma_i$ of [5000.0, 2000.0] and $\tau_i$ = [0.1, 1.0] ps.
    We shift time by dt to visualize $\Gamma(0)$.}
\end{figure}

\subsection{MD data Ala9}
For the MD data of ala9 at different discretizations, we compare the two methods.
\begin{figure}[H]
    \centering
    \includegraphics{/net/data02/greedisgod/paper1/figures/halfstep_ala9_md.pdf}
    \caption{We extract a memory kernel from equations 3 and 4 in the ala9 potential
    from the 10 $\mu$s data of ala9. We subsample the data at different time steps
    but keep the free energy constant for the memory extractions.
    We shift time by 1e-4 ps to visualize $\Gamma(0)$.}
\end{figure}

\section{Discussion}
The half-step method could only produce a better result in the case of no potential and a delta kernel.
In all cases, both methods were able to extract the delta kernel. For the discrete data of Ala9, both methods find hardly any
delta kernel contribution. The delta-kernel estimated by the full-step method is slightly larger than the
half-step method.

\end{document}
